\chapter{Приложение к методам второго порядка для вариационных неравенств}\label{ch:viqa}

\section{Постановка и место в работе}\label{sec:viqa_setup}

Главы~\ref{ch:sp-broyden-theory} и~\ref{ch:symmetric} развивают
\emph{локальную} теорию: они количественно ограничивают ошибку
аппроксимации якобиана $\Norm{B_k - J(x_k)}$, но скорость самого
итерационного процесса гарантируется лишь в~окрестности решения. В
настоящей главе показано, куда эта ошибка входит в~\emph{глобальном}
методе второго порядка~--- том самом верхнем уровне, ради которого
количественная оценка $\Norm{B_k - J(x_k)}$ и~представляет интерес
(ср.~Введение и~Заключение). Изложение следует
работе~\cite{viqa2024}\footnote{Результаты настоящей главы получены
в~соавторстве.}: оно сжато до~интерфейса с~главами~\ref{ch:sp-broyden-theory}
и~\ref{ch:symmetric}, а~формулировки теорем приведены как внешний
контекст и~атрибутированы по~этому источнику.

Рассматривается вариационное неравенство (ВН) Минти: найти
$x^*\in\mathcal{X}$ такое, что
\begin{equation}\label{eq:mvi}
  \langle F(x),\, x - x^*\rangle \ge 0 \qquad \forall\, x\in\mathcal{X},
\end{equation}
где оператор $F\colon\mathcal{X}\to\RR^n$ предполагается
\emph{монотонным}, $\langle F(x)-F(y),\,x-y\rangle\ge 0$, и~$L_1$-гладким,
$\Norm{\nabla F(x)-\nabla F(y)}_{\mathrm{op}}\le L_1\Norm{x-y}$.
Класс~\eqref{eq:mvi} охватывает как минимизацию ($F=\nabla f$), так и
минимаксные задачи $\min_x\max_y\Phi(x,y)$ с~оператором
$F(z)=(\nabla_x\Phi,\,-\nabla_y\Phi)$,~--- именно тот несимметричный
источник $F$, который выделен в~таблице~\ref{tab:ch1_classes}
главы~\ref{ch:jacobian}. Методы второго порядка для~\eqref{eq:mvi}
сходятся быстрее методов первого порядка, но~требуют на~каждом шаге
якобиан $\nabla F$, вычисление которого дорого; замена точного якобиана
аппроксимацией $J(v)\approx\nabla F(v)$ и~есть та неточность, цена
которой анализируется ниже.

\section{Неточный якобиан и оптимальная скорость}\label{sec:viqa_rate}

Метод VIJI~\cite{viqa2024} строится по~схеме дуальной экстраполяции и
на~каждой итерации решает вспомогательное ВН с~линейной по~$x$ моделью
оператора
\begin{equation}\label{eq:viqa_model}
  \Psi_v(x) = F(v) + J(v)\,[x-v],
\end{equation}
в~которой $J(v)$~--- приближение якобиана $\nabla F(v)$, играющее ту же
роль, что и~матрица $B_k$ в~главах~\ref{ch:sp-broyden-theory}
и~\ref{ch:symmetric}. Неточность модели~\eqref{eq:viqa_model}
контролируется одним числом:
\begin{equation}\label{eq:viqa_delta}
  \Norm{\nabla F(v) - J(v)}_{\mathrm{op}} \le \delta .
\end{equation}
Центральный результат~--- явная зависимость скорости от~уровня
неточности~$\delta$: для зазора ВН
$\mathrm{gap}(\widetilde{x}_T)$ после $T$~итераций~\cite{viqa2024}
\begin{equation}\label{eq:viqa_rate}
  \mathrm{gap}(\widetilde{x}_T)
  = O\!\left(\frac{L_1 D^3}{T^{3/2}} + \frac{\delta D^2}{T}\right),
\end{equation}
где $D$~--- диаметр $\mathcal{X}$. Первое слагаемое~--- оптимальная
скорость точного метода второго порядка; второе~--- штраф за~неточность
якобиана, линейный по~$\delta$. При выполнении более слабого
\emph{относительного} условия на~ошибку, $\Norm{(\nabla F(v)-J(v))[x-v]}
\le \tfrac{L_1}{2}\Norm{x-v}^2$, штрафное слагаемое поглощается и~метод
восстанавливает точную скорость $O(L_1 D^3/T^{3/2})$~\cite{viqa2024}.
Совпадающая нижняя оценка $\Omega(\delta D^2/T + L_1 D^3/T^{3/2})$,
также установленная в~\cite{viqa2024}, показывает, что
зависимость~\eqref{eq:viqa_rate} от~$\delta$ неулучшаема: уровень
неточности якобиана входит в~скорость по~существу, а~не~как артефакт
анализа. Тем самым любое количественное усиление оценки
$\Norm{B_k-J}$~--- предмет глав~\ref{ch:sp-broyden-theory}
и~\ref{ch:symmetric}~--- напрямую переносится в~оценку~\eqref{eq:viqa_rate}
внешнего метода.

\section{Квазиньютоновский ингредиент: связь с главами~\ref{ch:sp-broyden-theory} и~\ref{ch:symmetric}}\label{sec:viqa_qn}

Оценка~\eqref{eq:viqa_rate} превращает задачу «удешевить метод второго
порядка» в~задачу «удержать $\delta$ малым дёшево», что и~есть
квазиньютоновская постановка. В~\cite{viqa2024} приближение $J(v)$
строится \emph{односекущими} обновлениями Бройдена с~ограниченной
памятью~--- L-Broyden,
\begin{equation}\label{eq:viqa_lbroyden}
  J_{i+1} = J_i + \frac{(y_i - J_i s_i)\,s_i^\top}{s_i^\top s_i},
  \qquad i = 0,\dots,m-1,
\end{equation}
и~его демпфированным вариантом; пары $(s_i,y_i)$ берутся либо из~истории
оператора ($s_i=z_{i+1}-z_i$, $y_i=F(z_{i+1})-F(z_i)$), либо из~якобиан-векторных
произведений $y_i=\nabla F(x)\,s_i$ на~случайных направлениях. Для них
доказана оценка $\delta\le(m{+}2)\,L_0$ (L-Broyden) и~$\delta\le 2L_0$
(демпфированный вариант)~\cite{viqa2024}. Здесь $L_0$~--- константа Липшица
\emph{оператора}~$F$, то~есть равномерная оценка нормы якобиана
$\Norm{\nabla F(\cdot)}_{\mathrm{op}}\le L_0$ (её не~следует смешивать
с~константой Липшица якобиана~$L_1$ из~\S\ref{sec:viqa_setup}). Эти
оценки~--- априорный потолок: демпфирование удерживает
$\Norm{J}_{\mathrm{op}}\le L_0$, откуда
$\delta\le\Norm{\nabla F}_{\mathrm{op}}+\Norm{J}_{\mathrm{op}}\le 2L_0$
при~любом числе обновлений; качество же аппроксимации проявляется в~том,
что вблизи решения $\delta$ оказывается много меньше этого потолка. Существенно, что~\eqref{eq:viqa_lbroyden}~---
обновление ранга~1: на~каждом шаге сохраняется лишь текущее секущее
условие, то~есть тот самый ранг-1 проектор, расширение которого
до~ранга~$p{+}1$ составляет содержание главы~\ref{ch:sp-broyden-theory}.

Обе линии работы относятся к~одной и~той же величине. Уровень неточности
$\delta$ из~\eqref{eq:viqa_delta} есть в~точности ошибка аппроксимации
якобиана: при~$J(v)=B_k$ выполнено
$\delta=\Norm{B_k-\nabla F(v)}_{\mathrm{op}}\le\Fro{B_k-\nabla F(v)}$~---
та~самая величина $\Norm{B_k-J}$, которую анализируют
главы~\ref{ch:sp-broyden-theory} и~\ref{ch:symmetric}. Различается \emph{вид}
доступного над~ней контроля. Односекущая формула~\eqref{eq:viqa_lbroyden}
даёт равномерный априорный потолок $\delta\le 2L_0$, не~зависящий ни~от
размерности, ни~от~траектории. Мультисекущее обновление удерживает
на~каждом шаге $p{+}1$ секущих условий вместо одного~--- строго больше
информации об~ошибке (следствие~\ref{cor:subspace-alignment})~--- и~даёт
для~$\Norm{B_k-J}$ оценку ограниченной деградации (теорема~\ref{thm:bd}; для
симметричного случая~--- композитное обновление SS-$\mathcal{U}$,
Лемма~\ref{lem:bd_composite}) с~константой
$C_{\mathrm{PB}} = L\,(1{+}C_s/2)\sqrt{p+1}\,\kappa_p$, зависящей
от~обусловленности окна~$\kappa_p$, а~\emph{не}~от размерности~$n$. Это,
однако, оценка иной природы, чем глобальный потолок~$2L_0$:
теорема~\ref{thm:bd}~--- \emph{локальное} рекуррентное неравенство с~членом
деградации $C_{\mathrm{PB}}^2\rho_k^2$, действующее в~окрестности решения,
где $\rho_k$~--- радиус окна последних шагов. Строгое согласование этих
двух уровней~--- подстановка мультисекущей оценки в~роль~$\delta$ метода
VIJI с~выводом замкнутой скоростной оценки~--- вынесено в~резюме главы как
направление дальнейшей работы.

\paragraph{Резюме.}\label{par:viqa-summary}
Глава помещает локальную теорию глав~\ref{ch:sp-broyden-theory}
и~\ref{ch:symmetric} в~контекст глобального метода второго порядка для
вариационных неравенств~\cite{viqa2024}: уровень неточности якобиана
$\delta=\Norm{\nabla F - J}$ входит в~скорость~\eqref{eq:viqa_rate}
линейно и~неулучшаемо, а~его удержание возлагается на~квазиньютоновскую
аппроксимацию. В~\cite{viqa2024} эту роль играет односекущий L-Broyden
с~оценкой $\delta\le(m{+}2)L_0$; мультисекущее обновление
главы~\ref{ch:sp-broyden-theory} и~SS-$\mathcal{U}$ главы~\ref{ch:symmetric}
дают для~той же величины оценку с~геометрической константой $\kappa_p$
вместо размерности~$n$, сохраняя на~каждом шаге $p{+}1$ секущих условий.
Тесная подстановка мультисекущей оценки ограниченной деградации в~роль
$\delta$ метода VIJI~--- и~сопутствующее численное сравнение с~односекущим
L-Broyden~--- остаётся естественным направлением дальнейшей работы
(ср.~Заключение).
